% Code-Listings + eigene Commands
% Themes: Light (Default) und Dark - Umschalten via \codeDarkThemetrue
%
% Quellen für das Setup:
% - listings-Package: https://ctan.org/pkg/listings
% - VS Code Themes (Microsoft, MIT): github.com/microsoft/vscode/tree/main/extensions/theme-defaults/themes
% - TypeScript-Keywords aus der TS Language Spec

\definecolor{iuGreen}{RGB}{120, 190, 32}

% Theme-Toggle -> Default Light
\newif\ifcodeDarkTheme
\codeDarkThemefalse

% Light-Palette (VS Code Light+ leicht entsättigt für Print)
\definecolor{codeKeywordLight}{RGB}{0, 100, 170}
\definecolor{codeTypeLight}{RGB}{38, 127, 153}
\definecolor{codeStringLight}{RGB}{163, 21, 21}
\definecolor{codeCommentLight}{RGB}{100, 140, 80}
\definecolor{codeNumberLight}{RGB}{120, 120, 120}
\definecolor{codeFunctionLight}{RGB}{121, 94, 38}
\definecolor{codeIdentifierLight}{RGB}{0, 0, 0}
\definecolor{codeFrameLight}{RGB}{200, 200, 210}
\definecolor{codeBgTSLight}{RGB}{246, 248, 251}
\definecolor{codeBgJSLight}{RGB}{251, 250, 244}
\definecolor{codeBgJSONLight}{RGB}{248, 250, 246}
\definecolor{codeBgPyLight}{RGB}{247, 247, 250}
\definecolor{codeBgGenericLight}{RGB}{245, 245, 245}
\definecolor{codeHeaderBgLight}{RGB}{232, 236, 244}
\definecolor{codeHeaderTextLight}{RGB}{40, 60, 90}

% Dark-Palette (VS Code Dark+ dark_plus.json)
\definecolor{codeKeywordDark}{HTML}{569CD6}
\definecolor{codeTypeDark}{HTML}{4EC9B0}
\definecolor{codeStringDark}{HTML}{CE9178}
\definecolor{codeCommentDark}{HTML}{6A9955}
\definecolor{codeNumberDark}{HTML}{B5CEA8}
\definecolor{codeFunctionDark}{HTML}{DCDCAA}
\definecolor{codeIdentifierDark}{HTML}{D4D4D4}
\definecolor{codeFrameDark}{HTML}{3C3C3C}
\definecolor{codeBgTSDark}{HTML}{1E1E2E}
\definecolor{codeBgJSDark}{HTML}{1F1E1A}
\definecolor{codeBgJSONDark}{HTML}{1E1F1C}
\definecolor{codeBgPyDark}{HTML}{1F1F22}
\definecolor{codeBgGenericDark}{HTML}{1E1E1E}
\definecolor{codeHeaderBgDark}{HTML}{252526}
\definecolor{codeHeaderTextDark}{HTML}{9CDCFE}

% Aktive Theme-Farben
\ifcodeDarkTheme
  \colorlet{codeKeyword}{codeKeywordDark}
  \colorlet{codeType}{codeTypeDark}
  \colorlet{codeString}{codeStringDark}
  \colorlet{codeComment}{codeCommentDark}
  \colorlet{codeNumber}{codeNumberDark}
  \colorlet{codeFunction}{codeFunctionDark}
  \colorlet{codeIdentifier}{codeIdentifierDark}
  \colorlet{codeFrame}{codeFrameDark}
  \colorlet{codeBgTS}{codeBgTSDark}
  \colorlet{codeBgJS}{codeBgJSDark}
  \colorlet{codeBgJSON}{codeBgJSONDark}
  \colorlet{codeBgPy}{codeBgPyDark}
  \colorlet{codeBgGeneric}{codeBgGenericDark}
  \colorlet{codeHeaderBg}{codeHeaderBgDark}
  \colorlet{codeHeaderText}{codeHeaderTextDark}
\else
  \colorlet{codeKeyword}{codeKeywordLight}
  \colorlet{codeType}{codeTypeLight}
  \colorlet{codeString}{codeStringLight}
  \colorlet{codeComment}{codeCommentLight}
  \colorlet{codeNumber}{codeNumberLight}
  \colorlet{codeFunction}{codeFunctionLight}
  \colorlet{codeIdentifier}{codeIdentifierLight}
  \colorlet{codeFrame}{codeFrameLight}
  \colorlet{codeBgTS}{codeBgTSLight}
  \colorlet{codeBgJS}{codeBgJSLight}
  \colorlet{codeBgJSON}{codeBgJSONLight}
  \colorlet{codeBgPy}{codeBgPyLight}
  \colorlet{codeBgGeneric}{codeBgGenericLight}
  \colorlet{codeHeaderBg}{codeHeaderBgLight}
  \colorlet{codeHeaderText}{codeHeaderTextLight}
\fi

\graphicspath{{figures/}}
\pgfplotsset{compat=1.18}
\sisetup{locale=DE}
\renewcommand{\figurename}{Abb.}
\renewcommand{\tablename}{Tab.}

% TypeScript-Dialekt - listings kennt kein TS, Keywords aus der TS Language Spec
\lstdefinelanguage{TypeScript}{
  keywords={break, case, catch, class, const, continue, debugger, default, delete,
      do, else, enum, export, extends, false, finally, for, function, if,
      implements, import, in, instanceof, interface, let, new, null, package,
      private, protected, public, return, super, switch, this, throw, true,
      try, typeof, var, void, while, with, yield, async, await, from, as,
      readonly, abstract, namespace, type, satisfies, keyof, infer, declare, is},
  keywordstyle=\color{codeKeyword}\bfseries,
  ndkeywords={string, number, boolean, any, unknown, never, object, Promise, Array,
      Record, Partial, Pick, Omit, Map, Set, void, Buffer},
  ndkeywordstyle=\color{codeType}\bfseries,
  identifierstyle=\color{codeIdentifier},
  sensitive=true,
  comment=[l]{//},
  morecomment=[s]{/*}{*/},
  commentstyle=\color{codeComment}\itshape,
  stringstyle=\color{codeString},
  morestring=[b]',
  morestring=[b]",
  morestring=[b]`,
}

\lstdefinelanguage{JSONext}{
  keywords={true, false, null},
  keywordstyle=\color{codeKeyword}\bfseries,
  sensitive=true,
  string=[s]{"}{"},
  stringstyle=\color{codeString},
  comment=[l]{//},
  commentstyle=\color{codeComment}\itshape,
}

% Basis-Style für alle Code-Listings
\lstdefinestyle{codebase}{
basicstyle=\ttfamily\footnotesize\color{codeIdentifier},
backgroundcolor=\color{codeBgGeneric},
keywordstyle=\color{codeKeyword}\bfseries,
commentstyle=\color{codeComment}\itshape,
stringstyle=\color{codeString},
numberstyle=\tiny\color{codeNumber},
numbers=left,
numbersep=10pt,
breaklines=true,
breakatwhitespace=true,
columns=fullflexible,
keepspaces=true,
showstringspaces=false,
tabsize=2,
frame=single,
rulecolor=\color{codeFrame},
framerule=0.4pt,
framesep=4pt,
xleftmargin=14pt,
xrightmargin=4pt,
aboveskip=4pt,
belowskip=10pt,
captionpos=b,
upquote=true,
literate={ä}{{\"a}}1 {ö}{{\"o}}1 {ü}{{\"u}}1 {Ä}{{\"A}}1 {Ö}{{\"O}}1 {Ü}{{\"U}}1 {ß}{{\ss}}1 {…}{{...}}3
}

\lstdefinestyle{ts}{
  style=codebase,
  language=TypeScript,
  backgroundcolor=\color{codeBgTS},
}

\lstdefinestyle{js}{
  style=codebase,
  language=Java,
  backgroundcolor=\color{codeBgJS},
}

\lstdefinestyle{jsonext}{
  style=codebase,
  language=JSONext,
  backgroundcolor=\color{codeBgJSON},
}

\lstdefinestyle{pythonext}{
  style=codebase,
  language=Python,
  backgroundcolor=\color{codeBgPy},
}

\lstset{style=codebase}

% Pfad-Header über Code-Listing Nutzung zB: \codepath{backend/src/msl/core.ts}
\newcommand{\codepath}[1]{%
  \par\nopagebreak
  \noindent\colorbox{codeHeaderBg}{%
    \parbox{\dimexpr\linewidth-2\fboxsep\relax}{%
      \color{codeHeaderText}\ttfamily\footnotesize\strut #1\strut%
    }%
  }\par\nopagebreak\vspace{-2pt}%
}

\addbibresource{bibliography.bib}
\makeatletter
\renewcommand*{\apablx@ifrevnameappcomma}{\@secondoftwo}
\makeatother
\DefineBibliographyExtras{ngerman}{%
  \renewcommand*{\finalandcomma}{}%
}

\hypersetup{
  colorlinks=true,
  linkcolor=black,
  citecolor=black,
  urlcolor=codeKeyword,
  pdfauthor={\thesisAuthor},
  pdftitle={\thesisTitle},
  pdfsubject={Bachelorarbeit},
}

\newcommand{\todo}[1]{\textcolor{red}{\textbf{[TODO: #1]}}}
\newif\ifincludeListOfFigures
\includeListOfFiguresfalse
\newif\ifincludeListOfTables
\includeListOfTablesfalse
